\chapter{Normal Dağılım, Standart Puanlar ve Örnekleme Değişkenliği}
\label{ch:clt-se}
\index{normal dağılım}
\index{standart puan}
\index{örnekleme değişkenliği}

\begin{hedefler}
Bu bölümün sonunda okuyucu:
\begin{itemize}
  \item normal dağılımın merkez, yayılım, simetri ve alan özelliklerini açıklayabilecek,
  \item ham puanı $z$ ve T puanına dönüştürerek göreli konumu yorumlayabilecek,
  \item standart puan ile yüzdelik sırasını birbirinden ayırabilecek,
  \item bireysel puan dağılımı ile bir istatistiğin örnekleme dağılımını karşılaştırabilecek,
  \item standart sapma ve standart hatayı doğru bağlamda kullanabilecek,
  \item Merkezi Limit Teoreminin koşullarını ve söylemediği şeyleri açıklayabilecektir.
\end{itemize}
\end{hedefler}

\begin{hazirlik}
Bir öğrencinin puanının ortalamadan 10 puan yüksek olması tek başına yeterli
bilgi midir? Standart sapmanın neden gerekli olduğunu açıklayın.
\end{hazirlik}

\section{PDR vakası: Farklı ölçeklerdeki puanları karşılaştırmak}

Bir öğrencinin akademik kaygı ölçeğinden 72, psikolojik dayanıklılık ölçeğinden
60 puan aldığı düşünülmektedir. Kaygı puanının yüksek, dayanıklılık puanının
daha düşük görünmesi yalnız ham puanlara bakılarak söylenebilir mi? Ölçeklerin
puan aralıkları, norm grubundaki ortalamaları ve standart sapmaları bilinmeden
bu iki sayı doğrudan karşılaştırılamaz.

Kaygı ölçeğinde norm ortalaması 50 ve standart sapma 10; dayanıklılık ölçeğinde
ortalama 40 ve standart sapma 8 olsun. Öğrencinin göreli konumları
\[
z_{\text{kaygı}}=\frac{72-50}{10}=2{,}20,
\qquad
z_{\text{dayanıklılık}}=\frac{60-40}{8}=2{,}50
\]
olarak bulunur. Öğrenci her iki ölçekte de ortalamanın üzerindedir; norm grubu
içindeki göreli konumu dayanıklılıkta biraz daha yüksektir. Bununla birlikte
yüksek $z$ puanı tek başına klinik tanı veya müdahale gereksinimi göstermez.

Bu vaka bölümün ilk yarısında bireysel puanın konumunu, ikinci yarısında ise
bir öğrenci grubundan hesaplanan ortalamanın belirsizliğini anlamak için
kullanılacaktır.

\section{Normal dağılım}
\index{normal dağılım!özellikleri}
\index{yoğunluk fonksiyonu}

Normal dağılım $\mu$ ortalaması çevresinde simetrik, çan biçimli sürekli bir
dağılımdır. $\sigma$ küçüldükçe eğri daralır; büyüdükçe yayılır. Yaklaşık
normal bir dağılımda gözlemlerin yaklaşık yüzde 68'i $\mu\pm\sigma$, yüzde
95'i $\mu\pm2\sigma$ aralığındadır \cite{degroot2012,hogg2019}.

Normal dağılımın yoğunluk fonksiyonu
\[
f(x)=\frac{1}{\sigma\sqrt{2\pi}}
\exp\!\left[-\frac{(x-\mu)^2}{2\sigma^2}\right],
\qquad -\infty<x<\infty
\]
biçimindedir. Temel İstatistik düzeyinde formülün ezberlenmesinden çok şu
özellikler önemlidir:

\begin{itemize}
  \item Eğri $mu$ çevresinde simetriktir; ortalama, medyan ve mod çakışır.
  \item Eğri altındaki toplam alan 1'dir ve alanlar olasılıkları temsil eder.
  \item $mu$ konumu, $sigma$ yayılımı belirler.
  \item Eğrinin kuyrukları eksene yaklaşır, fakat kuramsal olarak eksene değmez.
  \item Belirli tek bir değerin olasılığı değil, bir aralığın alanı yorumlanır.
\end{itemize}

\begin{figure}[htbp]
\centering
\begin{tikzpicture}[alt={Normal dağılım eğrisi üzerinde ortalama ve bir ile iki standart sapmalık aralıklar},x=1cm,y=1cm]
  \draw[->,PrismGray] (-4.2,0) -- (4.4,0) node[right,font=\scriptsize]{$x$};
  \draw[PrismBlue,very thick,smooth,samples=100,domain=-4:4]
    plot (\x,{2.8*exp(-\x*\x/2)});
  \fill[PrismSubheadingBlue!18,smooth,samples=80,domain=-1:1]
    (-1,0) -- plot (\x,{2.8*exp(-\x*\x/2)}) -- (1,0) -- cycle;
  \draw[PrismBlue,dashed] (0,0) -- (0,2.85);
  \foreach \x/\lab in {-2/{\mu-2\sigma},-1/{\mu-\sigma},0/\mu,1/{\mu+\sigma},2/{\mu+2\sigma}}
    \draw (\x,0.08)--(\x,-0.08) node[below,font=\scriptsize]{$\lab$};
  \node[font=\scriptsize,text=PrismBlue] at (0,1.05) {yaklaşık \%68};
\end{tikzpicture}
\caption{Normal dağılımda merkez ve standart sapma birimleri.}
\end{figure}

\subsection{68--95--99,7 kuralı}
\index{68--95--99,7 kuralı}

Yaklaşık normal bir dağılımda gözlemlerin yaklaşık yüzde 68'i
$\mu\pm\sigma$, yüzde 95'i $\mu\pm2\sigma$ ve yüzde 99,7'si
$\mu\pm3\sigma$ aralığındadır. Ortalaması 50, standart sapması 10 olan bir
ölçekte puanların yaklaşık yüzde 68'i 40--60, yüzde 95'i 30--70 arasında
beklenir.

Bu kural yaklaşık bir betimleme sağlar. Gözlenen verinin belirgin çarpıklık,
çoklu tepe veya ağır kuyruk göstermesi durumunda normal eğriye dayalı bu
yüzdeler uygun olmayabilir.

\begin{etkinlik}
Ortalaması 40, standart sapması 6 olan yaklaşık normal bir puan dağılımında
$\mu\pm\sigma$, $\mu\pm2\sigma$ ve $\mu\pm3\sigma$ aralıklarını yazın.
Aralıkların her birinde yaklaşık kaç öğrencinin bulunacağını 1.000 kişilik bir
grup için tahmin edin.
\end{etkinlik}

\section{Normal eğri altında olasılık}
\index{normal dağılım!alan ve olasılık}

Sürekli dağılımda olasılık, iki değer arasında kalan alandır. Simetri
nedeniyle ortalamanın solunda ve sağında alanın yarısı bulunur:
\[
P(X<\mu)=P(X>\mu)=0{,}50.
\]
Standartlaştırma sonrasında bütün normal dağılımlar standart normal dağılım
tablosu veya yazılımıyla incelenebilir. Örneğin $z=1$ değerinin solundaki alan
yaklaşık 0,8413, sağındaki alan 0,1587'dir.

\begin{mistakebox}
Yoğunluk eğrisinin dikey yüksekliği doğrudan olasılık değildir. Olasılığı
eğrinin altında bir aralık boyunca kalan alan verir. Sürekli bir değişkende
tam olarak tek bir değeri gözleme olasılığı kuramsal olarak sıfırdır.
\end{mistakebox}

\section{Standart puan}
\index{z puanı@$z$ puanı}

Bir $x$ gözleminin ortalamadan kaç standart sapma uzakta olduğunu gösteren
standart puan
\[
z=\frac{x-\mu}{\sigma}
\]
biçimindedir. Örneklem özetleri kullanıldığında $z=(x-\bar{x})/s$ hesaplanır.
Pozitif $z$ ortalamanın üzerinde, negatif $z$ altında bir gözlemi gösterir.

$z$ puanının işareti yönü, mutlak değeri ortalamaya uzaklığı gösterir:

\begin{description}
  \item[$z=0$] Puan ortalamaya eşittir.
  \item[$z=1{,}5$] Puan ortalamanın 1,5 standart sapma üzerindedir.
  \item[$z=-0{,}8$] Puan ortalamanın 0,8 standart sapma altındadır.
\end{description}

Ham puana geri dönüş
\[
x=\mu+z\sigma
\]
ile yapılır. Örneğin $mu=50$, $sigma=10$ ve $z=-1{,}2$ için
$x=50-12=38$ bulunur.

\subsection{T puanı ve yüzdelik sırası}
\index{T puanı@$T$ puanı}
\index{yüzdelik sırası}
\index{norm grubu}

Negatif ve ondalıklı değerleri azaltmak için bazı ölçme uygulamalarında
\[
T=50+10z
\]
dönüşümü kullanılır. Kaygı örneğindeki $z=2{,}20$, $T=72$ değerine karşılık
gelir. Bu dönüşüm göreli konumu değiştirmez; yalnız ölçeğin başlangıç ve birim
değerini değiştirir.

Normal dağılım varsayımı altında $z=1{,}50$ yaklaşık 93,3. yüzdeliğe karşılık
gelir. Bu, puanın norm grubunun yaklaşık yüzde 93,3'ünden yüksek olduğu
anlamındadır. ``Yüzde 93,3 doğru yaptı'' veya ``özelliğin yüzde 93,3'üne
sahiptir'' anlamına gelmez.

Farklı ölçeklerdeki puanlar standartlaştırılarak karşılaştırılabilir; ancak
standartlaştırma ölçümlerin güvenilirlik veya içerik bakımından eşdeğer olduğu
anlamına gelmez.

\begin{researchbox}
Standart puan her zaman kullanılan norm grubuna görelidir. Bir üniversitenin
öğrencilerinden elde edilen ortalama ve standart sapmayla hesaplanan $z$ puanı,
ülke genelindeki öğrenciler için aynı göreli konumu göstermeyebilir. Norm
grubunun kimlerden, ne zaman ve hangi koşullarda oluştuğu raporlanmalıdır.
\end{researchbox}

\section{Dağılımın normalliğini değerlendirmek}
\index{normallik}
\index{Q--Q grafiği}
\index{Shapiro--Wilk testi}

Normal dağılım varsayımı tek bir karar kutusuna indirgenmemelidir. Histogram,
kutu grafiği ve normal Q--Q grafiği birlikte incelenir. Q--Q grafiğinde
noktaların yaklaşık bir doğru çevresinde ilerlemesi normal dağılımla uyumu;
sistematik eğrilik, ağır kuyruk veya belirgin uç noktalar sapmayı düşündürür.

Küçük örneklemde grafik düzensiz görünebilir; büyük örneklemde ise önemsiz bir
sapma normallik testinde anlamlı bulunabilir. Bu nedenle yalnız Shapiro--Wilk
gibi bir testin $p$-değerine bakarak yöntem değiştirmek dengeli bir yaklaşım
değildir. Araştırma tasarımı, örneklem hacmi, sapmanın biçimi, aykırı değerler
ve kullanılacak yöntemin sağlamlığı birlikte değerlendirilir.

\begin{mistakebox}
``Normallik testi anlamlı değilse veri normaldir'' sonucu kesin değildir.
Anlamlı olmayan sonuç, özellikle küçük örneklemde sapmayı gösterecek yeterli
kanıt bulunmadığını söyler; kusursuz normalliği kanıtlamaz.
\end{mistakebox}

\section{Bireysel dağılım ve örnekleme dağılımı}
\index{evren dağılımı}
\index{örneklem dağılımı}
\index{örnekleme dağılımı}
\index{standart hata}

Ham puanların dağılımı bireyler arasındaki değişkenliği gösterir. Aynı evrenden
tekrar tekrar $n$ birimlik örneklemler seçip her örneklemde $\bar X$
hesaplarsak, bu ortalamaların dağılımına $\bar X$'ın örnekleme dağılımı denir.
Bu iki dağılım birbirine karıştırılmamalıdır.

Üç düzeyi ayrı tutmak yararlıdır:

\begin{description}
  \item[Evren dağılımı] Hedef evrendeki bütün bireysel puanların dağılımıdır.
  \item[Örneklem dağılımı] Seçilen tek örneklemdeki bireysel puanların dağılımıdır.
  \item[Örnekleme dağılımı] Aynı büyüklükteki bütün olası örneklemlerden hesaplanan bir istatistiğin dağılımıdır.
\end{description}

Örnekleme dağılımını doğrudan gözlemlemeyiz; gerçekte çoğunlukla yalnız bir
örneklem seçilir. Kuramsal sonuçlar veya bilgisayar benzetimleri, istatistiğin
tekrarlı örneklemedeki davranışını anlamamızı sağlar.

Bağımsız ve aynı dağılımlı gözlemler için
\[
\E(\bar X)=\mu,
\qquad
\Var(\bar X)=\frac{\sigma^2}{n},
\qquad
SE(\bar X)=\frac{\sigma}{\sqrt n}.
\]
Uygulamada $\sigma$ bilinmiyorsa $SE(\bar X)\approx s/\sqrt n$ kullanılır.

Bu eşitlikler örneklem ortalamasının uzun dönemde evren ortalaması çevresinde
merkezlendiğini ve $n$ arttıkça dağılımının daraldığını gösterir. Bireysel
puanların standart sapması küçülmek zorunda değildir; daralan, ortalamaların
örnekleme dağılımıdır.

\begin{mistakebox}
Standart sapma bireysel gözlemlerin, standart hata ise bir tahmin edicinin
örneklemden örnekleme değişkenliğini ölçer. İkisi aynı nicelik değildir.
\end{mistakebox}

\section{Merkezi Limit Teoremi}
\index{Merkezi Limit Teoremi}

Sonlu varyans altında örneklem hacmi arttıkça standartlaştırılmış örneklem
ortalaması
\[
\frac{\bar X-\mu}{\sigma/\sqrt n}
\]
yaklaşık standart normal dağılıma yaklaşır. Bu sonuç, ham veri tam olarak
normal olmasa bile birçok güven aralığı ve test yönteminin neden büyük
örneklemlerde çalıştığını açıklar.

MLT'nin uygulanmasında üç nokta önemlidir:

\begin{enumerate}
  \item Gözlemler bağımsız veya bağımlılık analizde dikkate alınmış olmalıdır.
  \item Evren dağılımının varyansı sonlu olmalıdır.
  \item Gerekli örneklem hacmi dağılımın biçimine bağlıdır; evrensel bir ``$n=30$ yeterlidir'' kuralı yoktur.
\end{enumerate}


Evren yaklaşık simetrikse küçük örneklemlerde bile ortalama dağılımı normal
yaklaşıma uygun olabilir. Güçlü çarpıklık, seyrek fakat çok büyük değerler veya
bağımlı gözlemler daha fazla dikkat gerektirir.

MLT, ``her veri seti normal olur'' demez. Yaklaşık normallik kazanan nicelik,
uygun koşullar altında örneklem ortalamasının örnekleme dağılımıdır.
Teoremin koşulları ve yakınsama çerçevesi için bkz.\
\cite{billingsley1995,dasgupta2011}.

\begin{ornek}[Örneklem ortalaması için olasılık]
Bir PDR ölçeğinin evren ortalaması 50, standart sapması 12 olsun. Bağımsız
$n=36$ öğrencilik örneklemlerde
\[
E(\bar X)=50,
\qquad
SE(\bar X)=\frac{12}{\sqrt{36}}=2
\]
olur. Örneklem ortalamasının 54'ten büyük olma olasılığı için
\[
z=\frac{54-50}{2}=2
\]
hesaplanır. Standart normal dağılımda sağ kuyruk alanı yaklaşık 0,0228'dir.
Yani evren ortalaması gerçekten 50 ise, 36 kişilik örneklemlerin yaklaşık
yüzde 2,3'ünde ortalama 54'ten büyük olur.
\end{ornek}

\begin{figure}[htbp]
\centering
\begin{tikzpicture}[alt={Örneklem hacmi büyüdükçe örneklem ortalaması dağılımının daralmasını gösteren üç eğri},x=0.86cm,y=0.86cm]
  \draw[->,PrismGray] (-4.2,0) -- (4.4,0);
  \draw[PrismGray!65,thick,smooth,samples=100,domain=-4:4]
    plot (\x,{1.25*exp(-\x*\x/8)});
  \draw[PrismSubheadingBlue,thick,smooth,samples=100,domain=-3:3]
    plot (\x,{2.0*exp(-\x*\x/2)});
  \draw[PrismBlue,very thick,smooth,samples=100,domain=-2:2]
    plot (\x,{2.85*exp(-2*\x*\x)});
  \node[font=\scriptsize,text=PrismGray] at (3.25,1.15) {$n$ küçük};
  \node[font=\scriptsize,text=PrismSubheadingBlue] at (2.35,1.75) {$n$ orta};
  \node[font=\scriptsize,text=PrismBlue] at (1.25,2.75) {$n$ büyük};
\end{tikzpicture}
\caption{Örneklem hacmi arttıkça ortalamanın örnekleme dağılımının daralması.}
\end{figure}

\section{Karekök yasası}
\index{karekök yasası}
\index{örneklem hacmi!standart hata ilişkisi}

Standart hata $1/\sqrt n$ oranında azalır. Belirsizliği yarıya indirmek için
örneklem hacmini yaklaşık dört katına çıkarmak gerekir. Bu nedenle daha fazla
verinin getirisi vardır; fakat getiri doğrusal değildir.

\begin{table}[htbp]
\centering
\caption{$\sigma=12$ için örneklem hacmi ve standart hata ilişkisi.}
\begin{tabular}{@{}rrrr@{}}
\toprule
$n$ & 9 & 36 & 144 \\
\midrule
$SE(\bar X)=12/\sqrt n$ & 4 & 2 & 1 \\
\bottomrule
\end{tabular}
\end{table}

Örneklem hacmi 9'dan 36'ya, yani dört katına çıktığında standart hata 4'ten
2'ye düşer. Standart hatayı yeniden yarıya indirmek için örneklem hacmini
144'e çıkarmak gerekir. Bu ilişki örneklem planlamasında maliyet ile hassasiyet
arasındaki ödünleşimi gösterir.

\section{Python ile örnekleme dağılımı benzetimi}

Aşağıdaki kod, belirgin biçimde sağa çarpık bir evrenden farklı büyüklüklerde
örneklemler seçer ve örneklem ortalamalarını saklar:

\begin{lstlisting}
import numpy as np
import matplotlib.pyplot as plt

rng = np.random.default_rng(209)
tekrar = 5000

fig, ax = plt.subplots(1, 3, figsize=(9, 3))
for panel, n in zip(ax, [1, 10, 40]):
    ortalamalar = np.array([
        rng.exponential(scale=10, size=n).mean()
        for _ in range(tekrar)
    ])
    panel.hist(ortalamalar, bins=30, edgecolor="white")
    panel.set_title(f"n = {n}")
    panel.set_xlabel("örneklem ortalaması")

plt.tight_layout()
plt.show()
\end{lstlisting}

$n=1$ paneli evrenin çarpık biçimini yansıtır. $n=10$ ve özellikle $n=40$
için ortalamaların dağılımı daha simetrik olur ve daralır. Benzetim MLT'nin
iki sonucunu görünür kılar: biçim normale yaklaşırken standart hata
$1/\sqrt n$ oranında azalır.

\section{SPSS ile standart puan ve normallik incelemesi}

\begin{enumerate}
  \item Standart puan üretmek için \textbf{Analyze $>$ Descriptive Statistics $>$ Descriptives} yolunu izleyin.
  \item İlgili değişkeni seçip \textbf{Save standardized values as variables} seçeneğini işaretleyin.
  \item Dağılımı incelemek için \textbf{Analyze $>$ Descriptive Statistics $>$ Explore} menüsünü açın.
  \item \textbf{Plots} bölümünde histogram ve normality plots seçeneğini kullanın.
  \item Histogram, Q--Q grafiği, aykırı değerler ve örneklem hacmini birlikte yorumlayın.
\end{enumerate}

\begin{outputguidebox}
Çıktıda önce ortalama ve standart sapmanın kullanılan norm grubuna ait olup
olmadığını kontrol edin. Sonra gözlenen minimum--maksimum, çarpıklık, histogram
ve Q--Q grafiğini inceleyin. Normallik testinin $p$-değerini tek başına karar
kuralı yapmayın. Standartlaştırılmış değişkende ortalamanın yaklaşık 0,
standart sapmanın yaklaşık 1 olması hesaplamanın iç denetimidir.
\end{outputguidebox}

\section{Bütünleştirici çözümlü örnek}

Akademik kaygı ölçeğinde norm ortalaması 50, standart sapma 10'dur. Bir
öğrenci 65 puan almış; 64 kişilik bir sınıfın ortalaması 53, standart sapması
8 bulunmuştur.

\begin{enumerate}
  \item Öğrencinin norm grubuna göre standart puanı $(65-50)/10=1{,}50$'dir.
  \item Normal yaklaşım altında öğrenci yaklaşık 93,3. yüzdeliktedir.
  \item Sınıf ortalamasının tahmini standart hatası $8/\sqrt{64}=1$'dir.
  \item Bireysel puanın $z$ değeri öğrencinin göreli konumunu; standart hata ise sınıf ortalamasının örnekleme belirsizliğini anlatır.
\end{enumerate}

\begin{researchbox}
Bir öğrencinin yüksek standart puanı, ölçeğin geçerli ve güvenilir biçimde
uygulandığı, uygun norm grubunun kullanıldığı ve puanın görüşme bağlamında
değerlendirildiği varsayımıyla anlam kazanır. İstatistiksel sıra, bireyin
gereksinimlerini tek başına tanımlamaz.
\end{researchbox}

\section{Bilimsel raporlama örneği}

``Akademik kaygı puanları yaklaşık simetrik bir dağılım göstermiştir
($\bar{x}=53{,}0$, $SS=8{,}0$, $n=64$). Sınıf ortalamasının tahmini standart
hatası 1,0 puandır. Norm ortalaması 50 ve standart sapması 10 olan ölçekte 65
puan alan öğrencinin standart puanı $z=1{,}50$'dir. Bu değer norm grubuna göre
yüksek bir göreli konumu göstermekle birlikte klinik karar olarak
yorumlanmamıştır.''

\section{Literatür veri setiyle IBM SPSS laboratuvarı}
\index{Student Performance veri seti!z puanı}
\index{SPSS!standartlaştırılmış değer}
\index{z puanı!SPSS uygulaması}
\index{Merkezi Limit Teoremi!benzetim}
\index{örnekleme dağılımı!SPSS benzetimi}

Bu laboratuvarda \emph{Student Performance} matematik dosyasındaki yıl sonu
notu \texttt{G3} kullanılacaktır \cite{cortez2008data,cortezsilva2008}.
Uygulamanın ilk kısmı öğrencilerin G3 puanlarını örneklem ortalaması ve
standart sapmasına göre standartlaştırır. İkinci kısmı, gözlenen G3
dağılımından yerine koyarak 36 kişilik 5.000 örnek seçer ve örneklem
ortalamalarının dağılımını oluşturur.

\begin{researchbox}
Bu uygulamadaki z-puanları dışsal bir norm grubuna değil, gözlenen 395
öğrencinin ortalaması ve standart sapmasına dayanır. Dolayısıyla öğrencinin
bu veri dosyasındaki göreli konumunu gösterir; ulusal norm puanı veya klinik
eşik değildir.
\end{researchbox}

\subsection{Araştırma görevleri}

\begin{enumerate}
  \item G3 için standartlaştırılmış değişkeni üretip ortalama ve standart sapmasını denetleyin.
  \item Ham puanı 15 ve 5 olan öğrencilerin z-puanlarını yorumlayın.
  \item Ham G3 dağılımının normal olmak zorunda olmadığını grafiklerle gösterin.
  \item G3 dağılımından 36 kişilik 5.000 örnek seçerek örneklem ortalamalarını kaydedin.
  \item Örneklem ortalamalarının merkezi ve standart sapmasını kuramsal değerlerle karşılaştırın.
\end{enumerate}

\subsection{SPSS syntax: z-puanları}

IBM SPSS \textbf{Descriptives} yordamındaki \textbf{Save standardized values
as variables} seçeneği syntax'ta \texttt{/SAVE} alt komutuyla kullanılır.
G3 için oluşturulan değişken varsayılan olarak \texttt{ZG3} adını alır.

\begin{lstlisting}[language={}]
DATASET ACTIVATE StudentMath.

DESCRIPTIVES VARIABLES=G3
 /SAVE
 /STATISTICS=MEAN STDDEV MIN MAX.

DESCRIPTIVES VARIABLES=ZG3
 /STATISTICS=MEAN STDDEV MIN MAX.

EXAMINE VARIABLES=G3 ZG3
 /PLOT HISTOGRAM NPPLOT
 /STATISTICS DESCRIPTIVES
 /CINTERVAL 95
 /MISSING LISTWISE.
\end{lstlisting}

\subsection{Z-puanı çıktısının erişilebilir özeti}

\begin{table}[htbp]
\centering
\caption{G3 ve standartlaştırılmış G3 değişkeninin kontrol özeti.}
\begin{tabular}{@{}lrrrrr@{}}
\toprule
Değişken & $n$ & Ortalama & SS & En küçük & En büyük \\
\midrule
G3 & 395 & 10,42 & 4,58 & 0 & 20 \\
ZG3 & 395 & 0,00 & 1,00 & $-2,27$ & 2,09 \\
\bottomrule
\end{tabular}
\end{table}

Standartlaştırılmış değişkenin ortalamasının yuvarlama düzeyinde 0 ve standart
sapmasının 1 olması hesaplamanın iç denetimidir. Örneğin G3 puanı 15 olan bir
öğrenci için
\[
z=\frac{15-10{,}415}{4{,}581}\approx1{,}00,
\]
yani puan örneklem ortalamasının yaklaşık bir standart sapma üzerindedir.
G3 puanı 5 için $z\approx-1{,}18$ bulunur. Bu işaret ve büyüklük göreli konumu
gösterir; başarının nedeni veya öğrencinin gelecekteki performansı hakkında
tek başına açıklama sağlamaz.

\begin{mistakebox}
ZG3 değişkeninin ortalamasının 0 ve standart sapmasının 1 olması, G3
dağılımının normal olduğunu kanıtlamaz. Doğrusal standartlaştırma dağılımın
biçimini değiştirmez; G3'teki yığılmalar ve sınırlar ZG3'te de korunur.
\end{mistakebox}

\subsection{SPSS syntax: örneklem ortalamalarının dağılımı}

Aşağıdaki MATRIX programı G3 sütununu belleğe alır. Her tekrarda gözlenen 395
değer arasından yerine koyarak 36 seçim yapar, ortalamayı kaydeder ve sonuçları
ayrı bir SPSS veri dosyasına yazar. Dosya adı ve yolu gerektiğinde kullanıcı
tarafından değiştirilmelidir.

\begin{lstlisting}[language={}]
DATASET ACTIVATE StudentMath.
SET SEED=2092026.

MATRIX.
 GET g3 /VARIABLES=G3.
 COMPUTE N=NROW(g3).
 COMPUTE B=5000.
 COMPUTE m=36.
 COMPUTE ortalamalar=MAKE(B,1,0).

 LOOP i=1 TO B.
  COMPUTE toplam=0.
  LOOP j=1 TO m.
   COMPUTE k=TRUNC(RV.UNIFORM(0,1)*N)+1.
   COMPUTE toplam=toplam+g3(k,1).
  END LOOP.
  COMPUTE ortalamalar(i,1)=toplam/m.
 END LOOP.

 SAVE ortalamalar
  /OUTFILE='g3-orneklem-ortalamalari.sav'
  /VARIABLES=ortalama_G3.
END MATRIX.

GET FILE='g3-orneklem-ortalamalari.sav'.
DATASET NAME OrneklemOrtalamalari WINDOW=FRONT.

DESCRIPTIVES VARIABLES=ortalama_G3
 /STATISTICS=MEAN STDDEV MIN MAX.

EXAMINE VARIABLES=ortalama_G3
 /PLOT HISTOGRAM NPPLOT
 /STATISTICS DESCRIPTIVES
 /CINTERVAL 95.

DATASET ACTIVATE StudentMath.
\end{lstlisting}

\subsection{MLT çıktısının analizi}

G3'ün gözlenen ortalaması 10,42 ve standart sapması 4,58 olduğundan 36 kişilik
örneklem ortalamasının beklenen standart hatası
\[
SE(\bar X)=\frac{4{,}581}{\sqrt{36}}\approx0{,}764
\]
olur. SPSS benzetiminde 5.000 ortalamanın merkezi yaklaşık 10,42, standart
sapması ise yaklaşık 0,76 olmalıdır. Rassal seçim nedeniyle son basamaklarda
küçük farklılıklar beklenir.

Normal yaklaşım altında örneklem ortalamalarının yaklaşık yüzde 95'i
\[
10{,}415\pm1{,}96(0{,}764)
=[8{,}92;\ 11{,}91]
\]
aralığında beklenir. Bu aralık tek tek öğrencilerin yüzde 95'ini kapsayan bir
puan aralığı değildir; 36 kişilik tekrar örneklemlerin \emph{ortalamalarına}
ilişkindir.

\begin{table}[htbp]
\centering
\caption{MLT benzetimi için kuramsal ve SPSS çıktı denetimleri.}
\begin{tabular}{@{}p{0.34\textwidth}p{0.25\textwidth}p{0.25\textwidth}@{}}
\toprule
Denetim & Kuramsal hedef & Benzetimde beklenen \\
\midrule
Ortalamanın merkezi & $\mu_{\bar X}=10{,}42$ & 10,42'ye yakın \\
Ortalamaların standart sapması & $SE\approx0{,}764$ & 0,76'ya yakın \\
Dağılım biçimi & Yaklaşık normal & Histogram ve Q--Q yaklaşık doğrusal \\
Tekrar sayısı & --- & 5.000 geçerli ortalama \\
\bottomrule
\end{tabular}
\end{table}

Ham G3 dağılımında 0 puanında yığılma ve 0--20 sınırları bulunmasına rağmen,
36 bağımsız seçimin ortalamaları ham puanlardan daha düzgün ve daha dar bir
dağılım gösterir. MLT'nin anlattığı normal yaklaşım ham gözlemlere değil,
belirli koşullar altında örneklem ortalamasının örnekleme dağılımına aittir.

\begin{outputguidebox}
İki çıktıyı karıştırmayın: \textbf{ZG3} öğrencilerin standartlaştırılmış ham
puanlarını, \textbf{ortalama\_G3} ise tekrar örneklemlerin ortalamalarını
gösterir. İlkinde standart sapmanın 1 olması tanım gereğidir; ikincisinde
standart sapmanın yaklaşık 0,76 olması örneklem ortalamasının standart
hatasını gösterir.
\end{outputguidebox}

\subsection{Örnek bulgu paragrafı}

``Yıl sonu matematik notlarının ortalaması 10,42 ve standart sapması 4,58'dir
($n=395$). Örneklem içi standartlaştırmada puanı 15 olan bir öğrencinin
z-puanı yaklaşık 1,00'dır. Gözlenen dağılımdan yerine koyarak seçilen 36
kişilik 5.000 örneklemin ortalamaları yaklaşık 10,42 çevresinde ve 0,76
standart sapmayla dağılmıştır. Örneklem ortalamalarının histogramı ham G3
dağılımından daha düzgün görünmüştür. Bu sonuç ham notların normal olduğunu
değil, örneklem ortalamasının dağılımının MLT uyarınca normal biçime
yaklaştığını göstermektedir.''

\section{R ile çözümlü uygulama}
\label{sec:r-b05}

\textbf{Soru.} Ortalaması ve standart sapması 10 olan üstel bir evrenden
$n=1,5,30$ hacimlerinde bağımsız örneklemler üretin. Örneklem ortalamalarının
merkezini ve standart hatasını karşılaştırın.

\begin{lstlisting}[style=prismR]
set.seed(2026)
tekrar <- 10000
hacimler <- c(1, 5, 30)
sonuc <- data.frame(n = hacimler,
                    ortalama = NA_real_,
                    ampirik_se = NA_real_,
                    kuramsal_se = 10 / sqrt(hacimler))
for (sira in seq_along(hacimler)) {
  hacim <- hacimler[sira]
  ortalamalar <- replicate(tekrar,
    mean(rexp(hacim, rate = 1 / 10)))
  sonuc$ortalama[sira] <- mean(ortalamalar)
  sonuc$ampirik_se[sira] <- sd(ortalamalar)
  hist(ortalamalar, breaks = 35, probability = TRUE,
       main = paste("n =", hacim), xlab = "Ortalama")
  abline(v = 10, col = "blue", lwd = 2)
}
print(sonuc)
c(se_25 = 10 / sqrt(25), se_100 = 10 / sqrt(100))
\end{lstlisting}

\begin{outputguidebox}
\textbf{Çözüm ve kuramsal kontrol.} Üç ortalama da yaklaşık 10 çevresinde
olmalıdır. Kuramsal standart hatalar sırasıyla 10, 4,4721 ve 1,8257'dir;
benzetimdeki standart sapmalar bu değerlere yaklaşır, tam eşit olmak zorunda
değildir. $n=25$ ve $n=100$ için standart hatalar 2 ve 1'dir.
\end{outputguidebox}

\textbf{Yorum.} R'deki üstel dağılımın \texttt{rate} parametresi ortalamanın
tersidir. \texttt{rate = 10} yazmak farklı bir evren üretir. Ham evren
çarpık kalırken ortalamaların dağılımı örneklem hacmi arttıkça daha simetrik
ve dar hale gelir. $n=30$ her evren için otomatik normallik güvencesi değildir.

\textbf{Pekiştirme ve yanıt.} Standart hatayı yarıya indirmek için hacim
dört katına çıkarılır. Aynı tohumun R ve Python'da aynı rastgele sayıları
üretmesi beklenmez; kuramsal hedefleri karşılaştırın.
\section{Bölüm sonu çözümlü problemler}

\begin{enumerate}
  \item \textbf{Standart puan.} Kaygı puanlarının ortalaması 50, standart sapması 8'dir. 66 puan alan öğrencinin $z$ puanını ve konumunu yorumlayın.

  \textbf{Çözüm:} $z=(66-50)/8=2$'dir. Öğrencinin puanı ortalamanın iki standart sapma üzerindedir. Bu göreli konum tek başına klinik tanı değildir.

  \item \textbf{Normal olasılık.} $X\sim N(100,15^2)$ ise 115'in altında kalma olasılığını bulun.

  \textbf{Çözüm:} $z=(115-100)/15=1$ olur. Standart normal dağılımdan $P(Z<1)\approx0{,}8413$; dolayısıyla olasılık yaklaşık yüzde 84,1'dir.

  \item \textbf{Merkezî limit teoremi.} Evren standart sapması 12 olan bir puan için $n=36$ kişilik örneklem ortalamasının standart hatasını bulun. Evren hafif çarpıksa normal yaklaşımın neden yine kullanılabileceğini açıklayın.

  \textbf{Çözüm:} $SE(\bar X)=12/\sqrt{36}=2$'dir. Bağımsız gözlemler ve sonlu varyans altında örneklem ortalamasının dağılımı, ham puanlar tam normal olmasa da örneklem büyüdükçe normale yaklaşır. Güçlü aykırı değer veya bağımlılık varsa yalnız $n$ değerine güvenilmez.
\end{enumerate}

\bolumkaynaklari{\cite{degroot2012,billingsley1995,wasserman2004,hogg2019,cortez2008data,cortezsilva2008}}

\section*{Hatırla--Uygula--Yansıt}

\begin{enumerate}
\renewcommand{\labelenumi}{\textbf{\Alph{enumi}.}}
  \item \textbf{Hatırla:} Normal dağılımın beş temel özelliğini açıklayın.
  \item \textbf{Hesapla:} $\bar x=70$, $s=10$ olan dağılımda $x=85$ için $z$ ve T puanını bulun.
  \item \textbf{Ters dönüştür:} $\mu=40$, $\sigma=6$ ve $z=-1{,}5$ için ham puanı hesaplayın.
  \item \textbf{Yorumla:} $z=1{,}28$ değerinin yaklaşık 90. yüzdelik olması ne anlama gelir, ne anlama gelmez?
  \item \textbf{Ayır:} Standart sapma ile standart hatayı aynı PDR veri seti üzerinden karşılaştırın.
  \item \textbf{Hesapla:} $s=12$ için $n=36$ ve $n=144$ durumlarında standart hataları bulun.
  \item \textbf{Eleştir:} ``Örneklem 30'dan büyük olduğu için veriler normaldir'' ifadesini düzeltin.
  \item \textbf{Tanıla:} Histogram yaklaşık simetrik, Q--Q grafiğinde tek güçlü aykırı değer varsa hangi kontrolleri yaparsınız?
  \item \textbf{Uygula:} Python benzetiminde $n=5,20,80$ değerlerini kullanıp biçim ve yayılım değişimini karşılaştırın.
  \item \textbf{Yansıt:} PDR uygulamasında standart puanın neden tek başına tanı veya müdahale kararı veremeyeceğini tartışın.
\end{enumerate}

\bolumbaglantisi{chapters/pdr/03-normal-z-mkt.tex}